% ══════════════════════════════════════════════════════════════════════════════
%  CHAPTER 2: STAR FORMATION
% ══════════════════════════════════════════════════════════════════════════════
\chapter{Star Formation}

% ═════════════════════════════════════════════════════════════════════════════ 
\bigskip\noindent \bigskip\noindent \bigskip\noindent
\section{Sources of Star Formation}

\subsection{Molecular Clouds}
\begin{intuition}
Molecular clouds are dense regions of gas and dust in the interstellar medium where star formation occurs. They are primarily composed of molecular hydrogen (H$_2$) and can have masses ranging from a few solar masses to millions of solar masses. The high density and low temperature of molecular clouds allow them to collapse under their own gravity, leading to the formation of stars.
\end{intuition}

% ═════════════════════════════════════════════════════════════════════════════ 
\bigskip\noindent \bigskip\noindent \bigskip\noindent
\section{Analysis of Formation Criteria}
\subsection{Hydrostatic Equilibrium}
\begin{foundation}
Assume spherical symmetry and radial dependence only.
Neglect rotation, magnetic fields, and turbulence in the first model to isolate gravity-versus-pressure physics.
\end{foundation}

\begin{theory}
\begin{align}
    dm(r,t) &= 4\pi r^2\rho\,dr, \\
    \frac{\partial m}{\partial t}+v\frac{\partial m}{\partial r} &= 0.
\end{align}
\end{theory}

\begin{example}{}{ex:radialaccel}
Derive the radial acceleration $\ddot{r}$ of a fluid element at radius $r$ in a star, given the mass enclosed within $r$ is $m(r)$ and the density is $\rho(r)$.

\smallskip\textbf{Solution.}
The gravitational force per unit mass at radius $r$ is given by:
\begin{equation}
g = \frac{G m(r)}{r^2}.
\end{equation}
The pressure gradient force per unit mass is given by:
\begin{equation}
f_P = -\frac{1}{\rho} \frac{dP}{dr}.
\end{equation}
The total radial acceleration $\ddot{r}$ is the sum of these two forces:
\begin{equation}
\ddot{r} = f_P + g = -\frac{1}{\rho} \frac{dP}{dr} + \frac{G m(r)}{r^2}.
\end{equation}
\end{example}

\begin{example}{}{ex:radialaccel2}
Consider a star of mass $M$ and radius $R$ in hydrostatic equilibrium. If the star is suddenly compressed to half its radius, what is the radial acceleration $\ddot{r}$ at the surface?

\smallskip\textbf{Solution.}
The gravitational force per unit mass at the surface is given by:
\begin{equation}
g = \frac{GM}{R^2}.
\end{equation}
If the star is compressed to half its radius, the new radius is $R' = R/2$. The new gravitational force per unit mass at the surface is:
\begin{equation}
g' = \frac{GM}{(R/2)^2} = \frac{GM}{R^2/4} = 4\frac{GM}{R^2} = 4g.
\end{equation}
The radial acceleration $\ddot{r}$ at the surface is then:
\begin{equation}
\ddot{r} = g' - g = 3g = 3\frac{GM}{R^2}.
\end{equation}
\end{example}

\begin{theory}{Hydrostatic equilibrium}
At hydrostatic equilibrium, the inward gravitational force is balanced by the outward pressure gradient and $\ddot{r}=0$:
\begin{equation}
    \frac{dP}{dr} = -\frac{G m(r) \rho}{r^2}    
\end{equation}
\end{theory}


\begin{example}{}{ex:hydrostatic}
Consider a sphere of mass $M$ and radius $R$ with a uniform density $\rho$.
    \begin{enumerate}[label=(\alph*)]
        \item Write the pressure gradient $-\dfrac{dP}{dr}$ at distance $r$ from the center at hydrostatic equilibrium.
        \item Calculate the pressure $P(r)$ at distance $r$ from the center.
    \end{enumerate}

\smallskip\textbf{Solution.}
\begin{enumerate}[label=(\alph*)]
    \item The mass enclosed within radius $r$ is given by:
    \begin{equation}
        m(r) = \frac{4}{3}\pi r^3 \rho.
    \end{equation}
    The hydrostatic equilibrium equation is:
    \begin{equation}
        -\frac{dP}{dr} = \frac{G m(r) \rho}{r^2}.
    \end{equation}
    Substituting $m(r)$, we get:
    \begin{equation}
        -\frac{dP}{dr} = \frac{G \left(\frac{4}{3}\pi r^3 \rho\right) \rho}{r^2} = \frac{4\pi G}{3} r \rho^2.
    \end{equation}
    
    \item To find $P(r)$, we integrate the pressure gradient from the surface ($r=R$, where $P(R)=0$) to a point inside the sphere:
    \begin{align*}
        P(r) &= -\int_r^R -\frac{dP}{dr'} dr' = \int_r^R \frac{4\pi G}{3} r' \rho^2 dr' \\
        &= \frac{4\pi G}{3} \rho^2 \int_r^R r' dr' = \frac{4\pi G}{3} \rho^2 \left[\frac{r'^2}{2}\right]_r^R \\
        &= \frac{4\pi G}{3} \rho^2 \left(\frac{R^2 - r^2}{2}\right) = 2\pi G \rho^2 (R^2 - r^2).
    \end{align*}
\end{enumerate}
\end{example}
% ═════════════════════════════════════════════════════════════════════════════ 
\bigskip\noindent \bigskip\noindent \bigskip\noindent
\section{Virial Theorem and Jeans Criterion}
% ----------------------------------------------------------------------------------
\subsection{Energetics of Self-Gravitating Gas}
\begin{theory}
\begin{align}
    E_{\text{GR}} &= -\int_0^M \frac{Gm}{r}\,dm = -f_{\text{GR}}\frac{GM^2}{R}, \\
    \langle P\rangle &= -\frac{1}{3}\frac{E_{\text{GR}}}{V}, \\
    P &= \frac{1}{3}nm\langle v^2\rangle = \frac{2}{3}n\langle E_k\rangle.
\end{align}
For a star with a uniform density profile (i.e., $\rho = \text{constant}$), $f_{\text{GR}}\approx 3/5$.
\end{theory}
% ----------------------------------------------------------------------------------
\bigskip\noindent \bigskip\noindent \bigskip\noindent
\subsection{Jeans Criterion}
\begin{intuition}
The Jeans criterion determines whether a cloud of gas will collapse under its own gravity to form a star. It compares the gravitational energy of the cloud to its thermal energy. If the gravitational energy exceeds twice the thermal energy, the cloud is unstable and will collapse.
\end{intuition}
\bigskip\noindent \bigskip\noindent \bigskip\noindent
\subsection{Jeans Instability}
\begin{theory}
\begin{align}
    |E_{\text{GR}}| &> 2E_{\text{th}} \quad \text{(collapse)}, \\
    M_J &= \left(\frac{5k_B T}{G\mu m_H}\right)^{3/2}\left(\frac{3}{4\pi\rho}\right)^{1/2}, \\
    \lambda_J &= \sqrt{\frac{\pi c_s^2}{G\rho}}, \\
    M_J &\propto T^{3/2}\rho^{-1/2}.
\end{align}
\end{theory}

\begin{importantnote}
Lower temperature and higher density both reduce $M_J$.
This is why cooling strongly promotes fragmentation and cluster formation.
\end{importantnote}

\begin{example}{}{ex:jeans}
Calculate the Jeans mass for a molecular cloud with a temperature of $10\,\text{K}$ and a density of $10^{-20}\,\text{kg/m}^3$. Assume the mean molecular weight $\mu$ is 2.3 (typical for molecular clouds).   

\smallskip\textbf{Solution.}
First, we need to calculate the sound speed $c_s$ using the formula:
\begin{equation}
c_s = \sqrt{\frac{k_B T}{\mu m_H}},
\end{equation}
where $k_B$ is the Boltzmann constant ($1.38 \times 10^{-23}\,\text{J/K}$), $T$ is the temperature, $\mu$ is the mean molecular weight, and $m_H$ is the mass of a hydrogen atom ($1.67 \times 10^{-27}\,\text{kg}$).
Substituting the values, we get:
\begin{equation}
c_s = \sqrt{\frac{1.38 \times 10^{-23} \times 10}{2.3 \times 1.67 \times 10^{-27}}} \approx 0.19\,\text{km/s}.
\end{equation}
Next, we can calculate the Jeans mass $M_J$ using the formula:
\begin{equation}    
M_J = \left(\frac{5k_B T}{G\mu m_H}\right)^{3/2}\left(\frac{3}{4\pi\rho}\right)^{1/2},
\end{equation}
where $G$ is the gravitational constant ($6.67 \times 10^{-11}\,\text{N m}^2/\text{kg}^2$) and $\rho$ is the density. Substituting the values, we get:
\begin{align}
M_J &= \left(\frac{5 \times 1.38 \times 10^{-23} \times 10}{6.67 \times 10^{-11} \times 2.3 \times 1.67 \times 10^{-27}}\right)^{3/2}\left(\frac{3}{4\pi \times 10^{-20}}\right)^{1/2} \\
&\approx 1.0 \times 10^{30}\,\text{kg} \approx 0.5\,M_\odot.
\end{align}
\end{example}


\section{End of Formation: Protostars and Pre-Main Sequence Stars}
